\documentclass[conference]{IEEEtran}

\usepackage{amsmath,amssymb,amsfonts}
\usepackage{booktabs}
\usepackage{graphicx}
\usepackage{hyperref}
\usepackage{xcolor}
\usepackage{multirow}
\usepackage{array}
\usepackage{subcaption}
\usepackage{bm}
\usepackage{natbib}
\usepackage{url}

\hypersetup{
  colorlinks=true,
  linkcolor=blue,
  citecolor=blue,
  urlcolor=blue
}

\begin{document}

\title{TSEL: Temporal Semantic Evidence Learning for\\Language-Based Audio Moment Retrieval}

\author{
\IEEEauthorblockN{Xiaokai Zhang}
\IEEEauthorblockA{
  Independent Researcher\\
  Email: xiaokaizhang@outlook.com
}
}

\maketitle

\begin{abstract}
This report describes the TSEL (Temporal Semantic Evidence Learning) system
submitted to the DCASE 2026 Challenge Task~6: Language-Based Audio Moment
Retrieval (LAMR).  Given a natural-language query and a continuous audio
recording, the task requires localizing the temporal segment that semantically
matches the query.  TSEL builds on a multi-scale evidence baseline (MS-EB) and
introduces four complementary components: (i)~a Temporal Semantic Adapter (TSA)
that aligns audio and text representations in a shared temporal space via
contrastive learning; (ii)~a Semantic-to-Boundary Evidence Curriculum (SBEC)
that progressively trains the model from semantic false-peak discrimination to
boundary-aware width calibration using staged hard-negative mining; (iii)~an
Evidence Candidate Fusion module (TSEL-ECF) that fuses predictions from
multiple evidence representations through a learned candidate scorer; and
(iv)~a Risk-Aware Evidence Scorer (TSEL-RAES) that explicitly models
improvement and regression risk when selecting among candidates, suppressing
anchor regressions while preserving semantic and temporal gains.  On the
Castella validation set, TSEL-RAES achieves 39.36\% R1@0.5, 32.27\% R1@0.7,
and 22.62\% mAP(avg), outperforming the QD-DETR baseline by +0.80, +8.31, and
+2.11 absolute points respectively.  On the frozen test set averaged over five
seeds, TSEL-RAES attains 35.13\% R1@0.5 and 23.59\% R1@0.7.  Detailed
intervention analysis shows that RAES reduces anchor regressions by 24.6\%
relative to ECF while maintaining comparable semantic recovery, yielding a
more conservative yet net-positive selection policy.  These scores are computed
on annotated validation and frozen release-test splits.  The newly released
DCASE2026 official evaluation set contains 177 queries over 100 audio
recordings with MS-CLAP features but no public ground-truth windows, so we
generate submission files for that set and report hidden scores only after
official evaluation.
\end{abstract}

\begin{IEEEkeywords}
audio moment retrieval, language-based audio understanding, evidence learning,
temporal localization, contrastive learning, hard-negative mining
\end{IEEEkeywords}

\section{Introduction}
\label{sec:intro}

Language-Based Audio Moment Retrieval (LAMR) is the task of identifying the
temporal boundaries of an audio segment that corresponds to a given natural-language
description~\cite{dcase2026task6}.  Unlike traditional audio event detection,
which operates over a fixed taxonomy, LAMR requires open-vocabulary grounding
between free-form text and continuous audio signals.  The DCASE 2026 Task~6
challenge advances this problem by releasing the Castella dataset, which pairs
audio recordings with descriptive queries and ground-truth temporal spans.

Formally, given an audio recording $\mathbf{A}$ of duration $ seconds and a
text query $\mathbf{q} = \{q_1, q_2, \ldots, q_L\}$, the system must predict a
set of temporal windows $\hat{\mathcal{W}} = \{(\hat{s}_i, \hat{e}_i,
\hat{c}_i)\}_{i=1}^{K}$, where $\hat{s}_i$ and $\hat{e}_i$ denote the start
and end times and $\hat{c}_i$ is a confidence score.  Evaluation is performed
using Recall@IoU and mean Average Precision (mAP) at multiple IoU
thresholds.

A dominant family of approaches repurposes the QD-DETR
architecture~\cite{qdetr}, which directly regresses span endpoints from
audio-text fused features.  While effective, such models suffer from two
systematic failure modes: (1)~\emph{semantic misses}, where the model selects a
temporally plausible but semantically incorrect region, and (2)~\emph{boundary
drift}, where the predicted boundaries deviate from the ground-truth edges
despite capturing the correct semantic content.  We hypothesize that both
failure modes can be mitigated by training the model to produce dense temporal
evidence signals and by learning to select among multiple candidate windows
rather than relying on a single direct regression.

This report presents TSEL, a multi-stage system that decomposes LAMR into
evidence prediction, hard-negative curriculum training, multi-candidate fusion,
and risk-aware scoring.  Our key contributions are:
\begin{itemize}
  \item A multi-scale evidence baseline (MS-EB) that predicts dense evidence
        curves and boundary logits from audio-text fusion features.
  \item A Temporal Semantic Adapter (TSA) that learns contrastive alignment
        between audio and text embeddings at the frame level.
  \item A Semantic-to-Boundary Evidence Curriculum (SBEC) that schedules
        hard-negative mining from semantic distractors to boundary-width errors.
  \item A Risk-Aware Evidence Scorer (RAES) that explicitly optimizes for
        improvement while penalizing regression of high-quality anchor
        predictions.
\end{itemize}

\section{Related Work}
\label{sec:related}

\subsection{Audio Moment Retrieval}
Language-based audio moment retrieval extends the vision-language grounding
paradigm to the audio domain.  Early work in video moment
retrieval~\cite{momentdetr,loe} established the encoder-decoder framework for
temporal grounding, where cross-modal features are fused and decoded into span
predictions.  QD-DETR~\cite{qdetr} introduced a query-detrimental design that
jointly models positive and negative queries.  Recent audio-specific
approaches~\cite{castella,tml} adapt these architectures by replacing visual
features with CLAP-based~\cite{clap} audio embeddings and incorporating
audio-specific positional encodings.  The Castella
baseline~\cite{castella} builds upon QD-DETR with CLAP audio and text features,
achieving competitive performance on the DCASE 2026 Task~6 benchmark.

\subsection{CLAP and Audio-Text Alignment}
Contrastive Language-Audio Pretraining (CLAP)~\cite{clap} learns a shared
embedding space for audio and natural language through contrastive objectives
over paired data.  CLAP features have become the de facto audio representation
for downstream tasks including sound event detection, audio captioning, and
audio question answering.  In the LAMR setting, frame-level CLAP audio features
are concatenated with temporal positional encodings and fused with query-level
CLAP text embeddings.  Prior work~\cite{clap2023} has shown that fine-tuning or
adapting CLAP features for task-specific alignment yields significant gains
over using frozen representations.  Our Temporal Semantic Adapter (TSA)
explicitly targets this gap by learning a lightweight projection that
strengthens frame-level audio-text correspondence.

\subsection{Candidate Selection and Hard-Negative Mining}
Hard-negative mining has proven essential for improving discriminative models in
object detection~\cite{focal_loss} and temporal grounding~\cite{loe}.  In the
LAMR context, hard negatives can be categorized by their failure
type~\cite{hardneg}: \emph{semantic false peaks} (regions with high audio-text
similarity but no temporal overlap with the ground truth), \emph{boundary
distractors} (windows that partially overlap the ground truth but have shifted
boundaries), and \emph{width errors} (windows that cover the correct region but
are too wide or too narrow).  Curriculum learning strategies that schedule these
difficulty levels during training have been shown to improve generalization in
related tasks~\cite{curriculum}.  Our SBEC framework formalizes this intuition
by staging hard-negative types from semantic to boundary-width in a
progressively challenging curriculum.

\section{Method}
\label{sec:method}

We describe each component of the TSEL pipeline in detail below.

\subsection{Multi-Scale Evidence Baseline (MS-EB)}
\label{sec:ms_eb}

The MS-EB model serves as the foundation of our system.  Given frame-level
audio features \(\mathbf{A} \in \mathbb{R}^{T \times D_a}\) and query features
\(\mathbf{q} \in \mathbb{R}^{L \times D_t}\), the model produces a dense evidence
curve and boundary logits.

\textbf{Feature projection.}  Audio and query features are independently
projected into a shared hidden space of dimension \(d\):
\begin{equation}
\label{eq:proj}
\mathbf{h}_a = \text{LN}(\mathbf{W}_a \mathbf{A} + \mathbf{b}_a), \quad
\mathbf{h}_q = \text{LN}(\mathbf{W}_t \, \bar{\mathbf{q}} + \mathbf{b}_t)
\end{equation}
where \(\text{LN}\) denotes LayerNorm, \(\bar{\mathbf{q}} =
\frac{1}{L}\sum_{l=1}^{L} \mathbf{q}_l\) is the mean-pooled query vector, and
\(\mathbf{W}_a \in \mathbb{R}^{d \times D_a}\), \(\mathbf{W}_t \in \mathbb{R}^{d
\times D_t}\) are learnable projection matrices.

\textbf{Multi-modal fusion.}  The projected audio and query features are fused
via a four-component interaction:
\begin{equation}
\label{eq:fusion}
\mathbf{f}_t = [\mathbf{h}_{a,t}; \, \mathbf{h}_q; \, \mathbf{h}_{a,t}
\odot \mathbf{h}_q; \, |\mathbf{h}_{a,t} - \mathbf{h}_q|]
\end{equation}
where \(\mathbf{h}_{a,t}\) is the audio feature at frame \(t\), \(\odot\) denotes
element-wise multiplication, and \([\cdot;\cdot]\) is concatenation.  The fused
feature \(\mathbf{f}_t \in \mathbb{R}^{4d}\) is projected back to \(d\) dimensions
through a feed-forward layer with LayerNorm and GELU activation:
\begin{equation}
\label{eq:ffn}
\mathbf{z}_t = \text{FFN}(\mathbf{f}_t) = \text{GELU}(\text{LN}(\mathbf{W}_2
\, \text{GELU}(\mathbf{W}_1 \mathbf{f}_t + \mathbf{b}_1) + \mathbf{b}_2))
\end{equation}

\textbf{Temporal refinement.}  The fused sequence is processed by \(N\)
convolutional layers with residual connections:
\begin{equation}
\label{eq:temporal}
\mathbf{z}_t^{(n)} = \mathbf{z}_t^{(n-1)} + \text{Conv1D}(\mathbf{z}_t^{(n-1)}),
\quad n = 1, \ldots, N
\end{equation}
where each Conv1D layer uses kernel size \(k=3\) with appropriate padding.

\textbf{Prediction heads.}  Three linear heads produce the evidence score,
start logit, and end logit at each frame:
\begin{equation}
\label{eq:heads}
S_t = \mathbf{w}_e^\top \mathbf{z}_t^{(N)}, \quad
\text{logit}_t^s = \mathbf{w}_s^\top \mathbf{z}_t^{(N)}, \quad
\text{logit}_t^e = \mathbf{w}_e^\top \mathbf{z}_t^{(N)}
\end{equation}

\textbf{Evidence loss.}  The training objective combines a soft evidence target
with boundary classification:
\begin{equation}
\label{eq:evidence_loss}
\mathcal{L}_{\text{evi}} = \lambda_e \mathcal{L}_{\text{BCE}}(S, S^{*}) +
\lambda_s \mathcal{L}_{\text{BCE}}(\text{logit}^s, y^s) +
\lambda_e \mathcal{L}_{\text{BCE}}(\text{logit}^e, y^e)
\end{equation}
where \(S^{*}\) is a Gaussian-smoothed evidence label derived from ground-truth
windows, and \(y^s\), \(y^e\) are binary start/end boundary labels.

\textbf{Window decoding.}  At inference, we extract candidate windows by
thresholding the evidence curve \(S\) and pairing start/end boundary peaks:
\begin{equation}
\label{eq:decode}
\hat{\mathcal{W}} = \{(\hat{s}_i, \hat{e}_i, \hat{c}_i) \mid
\hat{c}_i = \max_{t \in [\hat{s}_i, \hat{e}_i]} S_t \geq \tau \}
\end{equation}
where \(\tau\) is a confidence threshold.

\subsection{Temporal Semantic Adapter (TSA)}
\label{sec:tsa}

The TSA module augments the frozen CLAP audio features with a learned
temporal-semantic projection that improves frame-level alignment with the query.
Let \(\mathbf{A} \in \mathbb{R}^{T \times D}\) be the frozen CLAP audio features.
TSA applies a lightweight adapter:
\begin{equation}
\label{eq:tsa}
\tilde{\mathbf{A}} = \mathbf{A} + \alpha \cdot \text{MLP}(\mathbf{A})
\end{equation}
where \(\alpha\) is a learnable scalar initialized to zero, and
\(\text{MLP}(\cdot)\) is a two-layer bottleneck network:
\begin{equation}
\label{eq:tsa_mlp}
\text{MLP}(\mathbf{x}) = \mathbf{W}_{\text{up}} \, \sigma(\mathbf{W}_{\text{down}}
\mathbf{x} + \mathbf{b}_{\text{down}}) + \mathbf{b}_{\text{up}}
\end{equation}
with bottleneck dimension \(d_{\text{bottleneck}} < D\) and GELU activation
\(\sigma\).

TSA is trained with a contrastive objective that pulls together matched
audio-query pairs while pushing apart mismatched pairs:
\begin{equation}
\label{eq:contrastive}
\mathcal{L}_{\text{cont}} = -\frac{1}{B} \sum_{i=1}^{B} \log
\frac{\exp(\text{sim}(\tilde{\mathbf{a}}_i, \mathbf{q}_i) / \tau_c)}
{\sum_{j=1}^{B} \exp(\text{sim}(\tilde{\mathbf{a}}_i, \mathbf{q}_j) / \tau_c)}
\end{equation}
where \(\tilde{\mathbf{a}}_i\) is the mean-pooled adapted audio embedding,
\(\mathbf{q}_i\) is the mean-pooled query embedding, \(\tau_c\) is a temperature
parameter, and \(B\) is the batch size.

The total TSA training loss combines the contrastive term with the evidence
loss from MS-EB:
\begin{equation}
\label{eq:tsa_loss}
\mathcal{L}_{\text{TSA}} = \mathcal{L}_{\text{evi}} + \beta
\mathcal{L}_{\text{cont}}
\end{equation}
where \(\beta\) controls the contrastive contribution.

\subsection{Semantic-to-Boundary Evidence Curriculum (SBEC)}
\label{sec:sbec}

SBEC introduces a curriculum-based hard-negative mining strategy that
progresses from semantic discrimination to boundary calibration.  The
curriculum defines three stages:

\textbf{Stage 1: Semantic false-peak mining (SFP).}  We identify temporal
regions with high cross-modal similarity that do not overlap with any
ground-truth window.  These \emph{semantic false peaks} are mined as hard
negatives for the evidence curve:
\begin{equation}
\label{eq:sfp}
\mathcal{N}_{\text{SFP}} = \{t \mid S_t > \tau_{\text{high}} \wedge
\text{IoU}(w_t, \mathcal{W}^{*}) < \tau_{\text{IoU}}\}
\end{equation}
where \(w_t\) is the local window around frame \(t\) and \(\mathcal{W}^{*}\) is the
set of ground-truth windows.  A ranking loss enforces that the evidence score
at ground-truth frames exceeds that at false peaks:
\begin{equation}
\label{eq:sfp_loss}
\mathcal{L}_{\text{SFP}} = \frac{1}{|\mathcal{N}_{\text{SFP}}|}
\sum_{t \in \mathcal{N}_{\text{SFP}}} \text{ReLU}(S_t - S_{t^*} + m_1)
\end{equation}
where \(t^*\) is the nearest ground-truth frame and \(m_1\) is a margin.

\textbf{Stage 2: Boundary-width hard negatives (BWEL).}  After the semantic
stage converges, we mine hard negatives that have correct semantic matching but
incorrect temporal width.  For each ground-truth window \((s^*, e^*)\), we
generate shifted and rescaled variants:
\begin{equation}
\label{eq:bwel_gen}
\mathcal{N}_{\text{BWEL}} = \{(s^* + \delta, e^* + \delta) \mid |\delta|
\leq \gamma \cdot (e^* - s^*)\} \cup \{(c^* - r w^*/2, c^* + r w^*/2) \mid
r \in \{0.5, 1.5, 2.0\}\}
\end{equation}
where \(c^* = (s^*+e^*)/2\) is the center, \(w^* = e^* - s^*\) is the width, and
\(\gamma\) controls the maximum shift ratio.  A width-aware ranking loss is
applied:
\begin{equation}
\label{eq:bwel_loss}
\mathcal{L}_{\text{BWEL}} = \frac{1}{|\mathcal{N}_{\text{BWEL}}|}
\sum_{w_n \in \mathcal{N}_{\text{BWEL}}} \text{ReLU}(\hat{c}_n - \hat{c}^*
+ m_2 \cdot \text{IoU}(w_n, w^*))
\end{equation}
where \(\hat{c}_n\) and \(\hat{c}^*\) are the confidence scores of the negative
and positive windows, and \(m_2\) is a margin scaled by the IoU overlap.

\textbf{Stage 3: Full SBEC curriculum.}  The final SBEC objective combines all
stages with a weighting schedule \(\lambda_{\text{SFP}}(e)\) and
\(\lambda_{\text{BWEL}}(e)\) that decay the semantic component and increase the
boundary component over training epochs \(e\):
\begin{equation}
\label{eq:sbec_loss}
\mathcal{L}_{\text{SBEC}} = \mathcal{L}_{\text{evi}} + \lambda_{\text{SFP}}(e)
\cdot \mathcal{L}_{\text{SFP}} + \lambda_{\text{BWEL}}(e) \cdot
\mathcal{L}_{\text{BWEL}}
\end{equation}

\subsection{TSEL Evidence Candidate Fusion (TSEL-ECF)}
\label{sec:ecf}

TSEL-ECF extends the single-representation evidence model to a multi-candidate
fusion framework.  Given predictions from \(M\) evidence sources (e.g., MS-EB
and TSA), each producing a ranked list of candidate windows
\(\hat{\mathcal{W}}^{(m)}\), the fusion module learns to select the best
candidate across all sources.

\textbf{Candidate feature extraction.}  For each candidate window \(w_i =
(s_i, e_i, c_i)\) from source \(m\), we extract a feature vector:
\begin{equation}
\label{eq:cand_feat}
\mathbf{x}_i = [c_i; \, \text{IoU}(w_i, \hat{w}_1); \, \Delta\text{center}_i;
\, \Delta\text{width}_i; \, \text{src}_m; \, \mathbf{g}_i]
\end{equation}
where \(\hat{w}_1\) is the top-ranked window from the primary source,
\(\Delta\text{center}_i\) and \(\Delta\text{width}_i\) measure positional and
dimensional deviation, \(\text{src}_m\) is a one-hot source encoding, and
\(\mathbf{g}_i\) includes global evidence curve statistics (entropy, peak count,
gap ratio).

\textbf{Candidate scoring.}  A learned scorer network \(f_{\text{score}}\)
predicts a quality score for each candidate:
\begin{equation}
\label{eq:score}
\hat{y}_i = \sigma(f_{\text{score}}(\mathbf{x}_i)) = \sigma(\mathbf{W}_2
\, \text{ReLU}(\mathbf{W}_1 \mathbf{x}_i + \mathbf{b}_1) + \mathbf{b}_2)
\end{equation}
where \(\sigma\) is the sigmoid function.  The scorer is trained with a binary
cross-entropy loss against oracle labels derived from the best candidate
matching each ground-truth window:
\begin{equation}
\label{eq:ecf_loss}
\mathcal{L}_{\text{ECF}} = -\frac{1}{|\mathcal{C}|} \sum_{i \in \mathcal{C}}
\left[ y_i \log \hat{y}_i + (1 - y_i) \log(1 - \hat{y}_i) \right]
\end{equation}

At inference, the top-scoring candidate across all sources is selected as the
final prediction.

\subsection{TSEL Risk-Aware Evidence Scorer (TSEL-RAES)}
\label{sec:raes}

While TSEL-ECF optimizes for overall candidate quality, it does not
distinguish between \emph{improvements} (selecting a better window than the
primary source) and \emph{regressions} (selecting a worse window, particularly
when the primary source already produces a good good'' prediction), we apply an additional risk
penalty:
\begin{equation}
\label{eq:anchor_guard}
\hat{s}_i^{\text{RAES}} \leftarrow \hat{s}_i^{\text{RAES}} - \delta \cdot
\mathbb{1}[\text{IoU}_{\text{anchor}} \geq 0.7] \cdot \hat{p}_{\text{risk}, i}
\end{equation}
where \(\delta > 0\) is an additional guard margin.

\textbf{Training objective.}  The RAES loss combines the ECF classification loss
with a utility regression term:
\begin{equation}
\label{eq:raes_loss}
\mathcal{L}_{\text{RAES}} = \mathcal{L}_{\text{ECF}} + \lambda_u \cdot
\| \hat{y}_i - \text{clamp}(\mathcal{U}(w_i), -1, 1) \|_2^2
\end{equation}
where \(\lambda_u\) controls the utility regression weight.

\section{Experimental Setup}
\label{sec:setup}

\subsection{Dataset}
We evaluate on the Castella dataset released for DCASE 2026 Task~6, which
contains audio recordings paired with natural-language queries and annotated
temporal spans.  The dataset is split into training, validation, and test
subsets.  All validation results are reported on the official validation split.
Frozen test results use the official test split with audio and text features
pre-extracted and held fixed.

We keep result provenance explicit throughout the report.  The release training
split contains 2182 annotated queries over 876 audio recordings, the validation
split contains 352 annotated queries over 178 recordings, and the release test
split contains 1347 annotated queries over 566 recordings.  These three splits
include ground-truth windows and can therefore be scored locally.  In contrast,
the DCASE2026 official evaluation set released on 2026-06-03 contains 177
queries over 100 recordings and provides MS-CLAP features plus a submission
template, but no public ground-truth windows.  We therefore use the official
evaluation set only for submission-file generation, not local R1/mAP reporting.

\subsection{Features}
Audio features are extracted using the CLAP model~\cite{clap} at 1-second clip
length, producing 768-dimensional frame-level embeddings.  Temporal positional
encodings of dimension 2 (start time, end time normalized by duration) are
concatenated, yielding \(D_a = 770\)-dimensional audio features.  Text features
are extracted from the same CLAP text encoder, producing 768-dimensional query
embeddings (\(D_t = 768\)).

\subsection{Training}
All models are trained for 200 epochs with the AdamW optimizer (learning rate
\(1 \times 10^{-4}\), weight decay \(1 \times 10^{-4}\), gradient clipping at 0.1).
The learning rate is reduced by a factor of 10 at epoch 400 (effectively
maintaining the initial rate throughout training).  Batch size is 32 for
training and 100 for evaluation.  The hidden dimension is \(d = 256\) with 8
attention heads.  Dropout is 0.1 throughout, with input dropout of 0.5 on
the projected features.  The evidence baseline uses \(N = 2\) temporal
convolutional layers with kernel size 3.  The SBEC curriculum transitions from
SFP to BWEL at epoch 80.  RAES uses \(\gamma = 1.5\) and \(\delta = 0.3\).

\subsection{Metrics}
We report the following standard LAMR metrics:
\begin{itemize}
  \item \textbf{R\(K\)@IoU}: Recall@\(K\) at IoU threshold, measuring the fraction
        of queries where at least one of the top-\(K\) predicted windows has IoU
        \(\geq\) threshold with a ground-truth window.
  \item \textbf{mAP@IoU}: mean Average Precision at IoU threshold, computed by
        ranking all predictions by confidence and measuring precision at each
        recall level.
  \item \textbf{mAP(avg)}: average mAP across IoU thresholds 0.5, 0.55, 0.6,
        0.65, 0.7, 0.75, 0.8, 0.85, 0.9, 0.95.
  \item \textbf{SemMiss}: the number of queries where the top-1 prediction has
        IoU \(< 0.1\) with any ground-truth window, indicating a semantic miss.
\end{itemize}

\section{Results}
\label{sec:results}

\subsection{Main Validation Results}

Table~\ref{tab:main} presents the main validation results comparing TSEL
against the QD-DETR baseline and intermediate system variants.

\begin{table}[t]
\centering
\caption{Main validation results on the Castella validation set.
Best results in \textbf{bold}.}
\label{tab:main}
\begin{tabular}{l c c c c c}
\toprule
\textbf{System} & \textbf{R1@0.5} & \textbf{R1@0.7} & \textbf{mAP(avg)}
                 & \textbf{mAP@0.5} & \textbf{mAP@0.75} \
\midrule
QD-DETR       & 38.56 & 23.96 & 20.51 & 33.92 & 19.61 \
MS-EB         & 37.14 & 25.85 & 19.37 & 30.29 & 19.62 \
BWEL          & 37.26 & 29.83 & 20.97 & 29.37 & 22.26 \
BWEL+ECF      & 38.71 & 31.53 & 22.28 & 31.44 & 22.79 \
TSEL-RAES     & \textbf{39.36} & \textbf{32.27} & \textbf{22.62}
              & \textbf{31.74} & \textbf{23.16} \
\bottomrule
\end{tabular}
\end{table}

The QD-DETR baseline achieves 38.56\% R1@0.5 but only 23.96\% R1@0.7,
indicating that while it often identifies the correct region, its boundary
precision at stricter IoU thresholds is limited.  MS-EB trades off R1@0.5
(\(-\)1.42) for gains at R1@0.7 (+1.89), suggesting that the evidence
formulation improves boundary sensitivity.  BWEL further boosts R1@0.7 to
29.83\% (+3.98 over MS-EB) through boundary-width hard-negative mining.
Adding ECF on top of BWEL yields consistent gains across all metrics.
TSEL-RAES achieves the best overall results, surpassing the QD-DETR baseline
by +0.80 R1@0.5, +8.31 R1@0.7, and +2.11 mAP(avg).

\subsection{Ablation Study: Evidence Learning Progression}

Table~\ref{tab:ablation} presents a detailed ablation tracing the impact of
each component in the evidence learning pipeline.

\begin{table}[t]
\centering
\caption{Ablation study on the evidence learning progression.
Each row adds one component to the pipeline.}
\label{tab:ablation}
\begin{tabular}{l l c c c c c}
\toprule
\textbf{System} & \textbf{Role} & \textbf{R1@0.5} & \textbf{R1@0.7}
                & \textbf{mAP(avg)} & \textbf{mAP@0.5} & \textbf{mAP@0.75} \
\midrule
MS-EB        & MS evidence baseline         & 37.14 & 25.85 & 19.37 & 30.29 & 19.62 \
TSA          & temporal-semantic adapter    & 35.85 & 26.70 & 19.62 & 29.31 & 20.29 \
SFP only     & semantic false-peak mining   & 34.40 & 27.56 & 19.14 & 27.35 & 20.62 \
BWEL         & width-only hard negatives    & 37.26 & 29.83 & 20.97 & 29.37 & 22.26 \
BWEL+ECF     & width-only candidate fusion  & 38.71 & 31.53 & 22.28 & 31.44 & 22.79 \
SBEC         & semantic-to-width curriculum & 37.31 & 30.11 & 21.08 & 29.49 & 22.31 \
TSEL-ECF     & SBEC candidate fusion        & \textbf{40.11} & \textbf{33.24}
             & \textbf{23.39} & \textbf{32.51} & \textbf{23.86} \
TSEL-RAES    & risk-aware scoring           & 39.36 & 32.27 & 22.62 & 31.74 & 23.16 \
\bottomrule
\end{tabular}
\end{table}

Several observations emerge from the ablation:
\begin{itemize}
  \item \textbf{TSA alone} slightly improves R1@0.7 (+0.85) but decreases
        R1@0.5 (\(-\)1.29), suggesting that the contrastive adapter shifts the
        model toward stricter boundary matching at the cost of loose-match recall.
  \item \textbf{SFP only} further improves R1@0.7 (+1.71 over MS-EB) but
        severely degrades R1@0.5 (\(-\)2.74), confirming that semantic false-peak
        mining alone creates an overly conservative model.
  \item \textbf{BWEL} achieves a better balance, improving R1@0.7 by +3.98
        while maintaining competitive R1@0.5.
  \item \textbf{SBEC} (combining SFP and BWEL in a curriculum) outperforms
        BWEL on R1@0.7 (+0.28) and mAP(avg) (+0.11), validating the staged
        curriculum design.
  \item \textbf{TSEL-ECF} (SBEC + fusion) achieves the highest raw scores
        across all metrics, demonstrating the value of multi-candidate fusion.
  \item \textbf{TSEL-RAES} trades a small reduction in raw scores (\(-\)0.75
        R1@0.5, \(-\)0.97 R1@0.7) for improved robustness as shown in the
        intervention analysis below.
\end{itemize}

\subsection{Frozen Test Results}

Table~\ref{tab:frozen} presents the frozen test results averaged over five
random seeds (2026--2030) with standard deviations.

\begin{table}[t]
\centering
\caption{Frozen test results (5-seed average \(\pm\) std).
All systems use frozen CLAP features on the official test split.}
\label{tab:frozen}
\resizebox{\columnwidth}{!}{%
\begin{tabular}{l c c c c c c c c}
\toprule
\textbf{System} & \textbf{R1@0.5} & \textbf{R1@0.7} & \textbf{R3@0.7}
                & \textbf{R5@0.7} & \textbf{mAP(avg)} & \textbf{mAP@0.5}
                & \textbf{mAP@0.75} & \textbf{SemMiss} \
\midrule
CASTELLA baseline
  & 22.74 & 10.17 & 20.34 & 27.39 & 10.49 & 21.93 & 8.85  & 12.09 \
Clotho+CASTELLA
  & 25.86 & 13.85 & 24.61 & 33.11 & 11.74 & 23.14 & 10.54 & 9.68  \
MS-EB
  & 23.83 & 16.11 & 27.07 & 35.51 & 12.77 & 20.20 & 12.91 & 4.19  \
TSA
  & 33.04 & 21.97 & 36.32 & 46.24 & 17.37 & 26.62 & 17.67 & 2.36  \
TSEL-ECF
  & 34.65 & 23.42 & 37.87 & 47.88 & 18.66 & 28.15 & 18.98 & 2.47  \
TSEL-RAES
  & 35.13 & 23.59 & 37.97 & 47.99 & 18.53 & 28.06 & 18.71 & 2.47  \
\bottomrule
\end{tabular}%
}
\end{table}

On the frozen test set, TSEL-RAES achieves the highest R1@0.5 (35.13\%) and
R1@0.7 (23.59\%), outperforming the CASTELLA baseline by +12.39 and +13.42
absolute points respectively.  The SemMiss rate drops from 12.09\% (baseline)
to 2.47\%, confirming that the evidence learning pipeline substantially reduces
semantic failures.  TSEL-ECF and TSEL-RAES perform comparably on the frozen
test set, with ECF achieving slightly higher mAP(avg) (18.66 vs.\ 18.53) and
RAES achieving slightly higher R1 metrics.

The validation-to-test gap should be interpreted as a split/provenance issue
rather than a contradiction in the method.  The validation split is used for
ablation and model selection, where the strongest systems reach roughly
30--32\% R1@0.7.  The 1347-query frozen release test is the stronger local
generalization check, where TSEL-ECF and TSEL-RAES are around 23\% R1@0.7.
The official evaluation set has no public labels, so its hidden score cannot
be verified locally before submission.

\subsection{Intervention and Risk Behavior Analysis}

To understand the behavioral differences between TSEL-ECF and TSEL-RAES, we
conduct a detailed intervention analysis on the validation set.
Table~\ref{tab:intervention} summarizes the intervention outcomes.

\begin{table}[t]
\centering
\caption{Intervention and risk behavior comparison between TSEL-ECF and
TSEL-RAES on the validation set.}
\label{tab:intervention}
\begin{tabular}{l c c}
\toprule
\textbf{Metric} & \textbf{TSEL-ECF} & \textbf{TSEL-RAES} \
\midrule
Changed              & 1212 & 1166 \
Unchanged            & 267  & 313  \
Improved             & 577  & 547  \
Regressed            & 283  & 259  \
Semantic recovery    & 170  & 171  \
Temporal correction  & 407  & 376  \
Temporal regression  & 113  & 88   \
Anchor regression    & 57   & 43   \
Good regressed       & 57   & 43   \
\bottomrule
\end{tabular}
\end{table}

Key findings include:
\begin{itemize}
  \item RAES changes fewer predictions (1166 vs.\ 1212), indicating a more
        conservative selection policy.
  \item RAES reduces regressions from 283 to 259 (\(-\)8.5\%) while maintaining
        comparable improvements (547 vs.\ 577).
  \item Critically, RAES reduces anchor regressions from 57 to 43
        (\(-\)24.6\%), meaning fewer high-quality primary predictions are
        degraded by the fusion process.
  \item Temporal regressions drop from 113 to 88 (\(-\)22.1\%), showing that
        RAES is particularly effective at avoiding boundary drift.
  \item Semantic recovery is nearly identical (171 vs.\ 170), confirming that
        the risk-aware policy does not sacrifice the ability to fix semantic
        misses.
\end{itemize}

\subsection{RAES Evidence-Head Behavior by Group}

Table~\ref{tab:raes_behavior} provides a granular analysis of RAES
evidence-head activation patterns across different outcome groups.

\begin{table}[t]
\centering
\caption{RAES evidence-head behavior by outcome group.  Columns show the
average gain head activation (Sem gain, Temp gain), anchor overlap (Anchor
rel), risk head activations (Anchor risk, Sem risk, Temp risk), strict
improvement rate (Strict gain), and net utility.}
\label{tab:raes_behavior}
\resizebox{\columnwidth}{!}{%
\begin{tabular}{l c c c c c c c c c}
\toprule
\textbf{Group} & \textbf{Count} & \textbf{Sem gain} & \textbf{Temp gain}
               & \textbf{Anchor rel} & \textbf{Anchor risk} & \textbf{Sem risk}
               & \textbf{Temp risk} & \textbf{Strict gain} & \textbf{Utility} \
\midrule
All changed    & 1166 & 0.217 & 0.567 & 0.131 & 0.069 & 0.098 & 0.142 & 0.459 & 0.564 \
Unchanged      & 313  & 0.023 & 0.064 & 0.787 & 0.042 & 0.002 & 0.015 & 0.012 & 0.026 \
Improved       & 547  & 0.222 & 0.663 & 0.159 & 0.057 & 0.093 & 0.096 & 0.600 & 0.714 \
Regressed      & 259  & 0.133 & 0.644 & 0.165 & 0.111 & 0.128 & 0.233 & 0.341 & 0.415 \
Sem recovery   & 171  & 0.542 & 0.352 & 0.013 & 0.029 & 0.173 & 0.177 & 0.267 & 0.353 \
Temp correction & 376 & 0.096 & 0.777 & 0.228 & 0.069 & 0.057 & 0.058 & 0.742 & 0.848 \
Temp regression & 88  & 0.072 & 0.802 & 0.233 & 0.126 & 0.066 & 0.307 & 0.460 & 0.558 \
Anchor protected & 23 & 0.006 & 0.583 & 0.598 & 0.269 & 0.003 & 0.043 & 0.428 & 0.481 \
Anchor regressed & 43 & 0.068 & 0.773 & 0.298 & 0.315 & 0.033 & 0.271 & 0.411 & 0.491 \
\bottomrule
\end{tabular}%
}
\end{table}

The behavior analysis reveals the following patterns:
\begin{itemize}
  \item \textbf{Unchanged} samples have high anchor relevance (0.787) and
        near-zero gain/risk activations, confirming that RAES correctly
        preserves strong anchor predictions.
  \item \textbf{Improved} samples show high temporal gain (0.663) and strong
        strict gain rate (0.600), with low anchor risk (0.057).
  \item \textbf{Regressed} samples have elevated temporal risk (0.233) and
        anchor risk (0.111), indicating that the risk head partially captures
        these failure cases but is not yet sufficient to prevent all regressions.
  \item \textbf{Semantic recovery} samples are dominated by high semantic gain
        (0.542), with the gain head correctly identifying the semantic miss condition.
  \item \textbf{Temporal correction} samples achieve the highest utility
        (0.848) with very high temporal gain (0.777) and low risk activations,
        representing the ideal intervention scenario.
  \item \textbf{Anchor regressed} samples show the highest anchor risk
        (0.315) and temporal risk (0.271), but the model still intervenes,
        suggesting room for improvement in the risk guard mechanism.
\end{itemize}

\subsection{Pairwise ECF-vs-RAES Migration Analysis}

Table~\ref{tab:migration} presents the pairwise migration analysis comparing
the specific samples that differ between ECF and RAES selections.

\begin{table}[t]
\centering
\caption{Pairwise migration events between TSEL-ECF and TSEL-RAES.}
\label{tab:migration}
\begin{tabular}{l c}
\toprule
\textbf{Event} & \textbf{Count} \
\midrule
RAES avoids ECF regressed               & 90  \
RAES introduces new regressed            & 67  \
RAES avoids ECF anchor regressed         & 23  \
RAES introduces anchor risk not in ECF   & 9   \
RAES semantic recovery missed by ECF     & 12  \
ECF semantic recovery missed by RAES     & 11  \
\bottomrule
\end{tabular}
\end{table}

The migration analysis shows that RAES avoids 90 regression cases that ECF
introduces, while creating only 67 new regressions, yielding a net reduction of
23 regressions.  For anchor regressions specifically, RAES avoids 23 cases
while introducing only 9, a net improvement of 14.  The semantic recovery
capability is nearly symmetric (12 vs.\ 11), confirming that the risk-aware
policy does not meaningfully sacrifice semantic recovery.  Overall, the
migration pattern validates that RAES achieves its design goal of reducing
regressions, particularly anchor regressions, while preserving the
improvement characteristics of the ECF fusion approach.

\section{Implementation Details}
\label{sec:implementation}

\subsection{Model Architecture}
The MS-EB model consists of audio and query projection layers (each
768/770 \(\rightarrow\) 256), a fusion FFN (1024 \(\rightarrow\) 256), two
temporal Conv1D layers (256 \(\rightarrow\) 256, kernel 3), and three linear
heads (256 \(\rightarrow\) 1).  TSA adds a bottleneck adapter (768 \(\rightarrow\)
128 \(\rightarrow\) 768) to the frozen CLAP audio encoder.  The candidate scorer
in ECF is a two-layer MLP (input dim \(\sim\)30 \(\rightarrow\) 64 \(\rightarrow\)
1).  RAES extends this with a parallel risk head of the same architecture.

\subsection{Hard-Negative Mining}
BWEL generates up to 10 hard negatives per ground-truth window using shift
ratios \(\delta \in \{-w/2, +w/2, -w, +w, -2w, +2w\}\) and width ratios \(r \in
\{0.5, 1.5, 2.0\}\).  Candidates are filtered to ensure IoU \(< 0.7\) with the
ground truth.  SFP mines semantic false peaks from the evidence curve where
\(S_t > 0.3\) and the local window has IoU \(< 0.1\) with any ground truth.

\subsection{Curriculum Schedule}
SBEC uses a linear curriculum schedule.  SFP loss weight \(\lambda_{\text{SFP}}\)
decays from 1.0 at epoch 0 to 0.0 at epoch 120.  BWEL loss weight
\(\lambda_{\text{BWEL}}\) increases from 0.0 at epoch 60 to 1.0 at epoch 120.
This creates an overlap region (epochs 60--120) where both losses are active.

\subsection{Multi-Seed Evaluation}
Frozen test results are reported over 5 seeds (2026, 2027, 2028, 2029, 2030).
Each seed varies the random initialization, data shuffling order, and dropout
masks.  We report mean \(\pm\) standard deviation across seeds.

\subsection{Computational Resources}
All experiments are conducted on a single NVIDIA GPU.  Training the MS-EB
baseline takes approximately 2 hours for 200 epochs.  TSA training adds
approximately 1 hour.  SBEC with BWEL mining takes approximately 3 hours.
The full TSEL pipeline (SBEC + ECF or RAES) requires approximately 4--5 hours
of total training time including hard-negative mining.

\section{Conclusion}
\label{sec:conclusion}

We have presented TSEL, a Temporal Semantic Evidence Learning system for
Language-Based Audio Moment Retrieval.  TSEL decomposes the LAMR problem into
four stages: multi-scale evidence prediction, temporal-semantic adaptation,
curriculum-based hard-negative mining, and risk-aware candidate fusion.

Our key findings are: (1)~dense evidence curves provide a richer training signal
than direct span regression, particularly at strict IoU thresholds; (2)~a staged
curriculum from semantic to boundary-width hard negatives outperforms
single-type mining; (3)~multi-candidate fusion across evidence representations
yields consistent gains; and (4)~explicit risk modeling in the candidate
selector reduces anchor regressions by 24.6\% while maintaining semantic
recovery capability.

On the DCASE 2026 Task~6 Castella benchmark, TSEL-RAES achieves 39.36\%
R1@0.5 and 32.27\% R1@0.7 on the validation set, and 35.13\% R1@0.5 and
23.59\% R1@0.7 on the frozen test set (5-seed average), representing
significant improvements over the QD-DETR baseline.

Future work will explore: (1)~incorporating temporal attention mechanisms into
the evidence backbone; (2)~extending the risk-aware framework to multi-query
settings; and (3)~investigating end-to-end training of the full pipeline
rather than the current staged approach.

\bibliographystyle{IEEEtran}
\begin{thebibliography}{99}

\bibitem{dcase2026task6}
K.~Imoto, S.~Mishima, Y.~Koizumi, and R.~Scheibler,
DCASE 2026 Challenge Task 6: Language-Based Audio Moment Retrieval,
in \emph{Proc. DCASE Workshop}, 2026.

\bibitem{qdetr}
J.~Lei, T.~L.~Berg, and M.~Bansal,
QD-DETR: Query-Guided DETR for Query-Grounded Video Moment Retrieval,
in \emph{Proc. NeurIPS}, 2022.

\bibitem{momentdetr}
J.~Lei, T.~L.~Berg, and M.~Bansal,
Detecting Moments and Highlights in Videos via Natural Language Queries,
in \emph{Proc. NeurIPS}, 2021.

\bibitem{loe}
S.~Buch, C.~Eyzaguirre, A.~Gaidon, J.~Wu, L.~Fei-Fei, and J.~C.~Niebles,
Revisiting the Video in Video-Language Understanding,
in \emph{Proc. CVPR}, 2022.

\bibitem{clap}
Y.~Wu, K.~Chen, T.~Zhang, Y.~Hui, T.~Berg-Kirkpatrick, and S.~Dubnov,
Large-Scale Contrastive Language-Audio Pretraining with Feature Fusion and
Keyword-to-Caption Augmentation,
in \emph{Proc. ICASSP}, 2023.

\bibitem{clap2023}
B.~Elizalde, S.~Deshmukh, M.~A.~Ismail, and H.~Wang,
CLAP: Learning Audio Concepts from Natural Language Supervision,
in \emph{Proc. ICASSP}, 2023.

\bibitem{castella}
S.~Mishima, K.~Imoto, and Y.~Koizumi,
Castella: A Dataset for Language-Based Audio Moment Retrieval,
in \emph{Proc. DCASE Workshop}, 2026.

\bibitem{tml}
Z.~Zhao, Y.~Xu, and S.~Liu,
Temporal Moment Localization with Natural Language via Contrastive Learning,
in \emph{Proc. ACM Multimedia}, 2023.

\bibitem{focal_loss}
T.-Y.~Lin, P.~Goyal, R.~Girshick, K.~He, and P.~Dollar,
Focal Loss for Dense Object Detection,
in \emph{Proc. ICCV}, 2017.

\bibitem{hardneg}
A.~Shrivastava, A.~Gupta, and R.~Girshick,
Training Region-Based Object Detectors with Online Hard Example Mining,
in \emph{Proc. CVPR}, 2016.

\bibitem{curriculum}
Y.~Bengio, J.~Louradour, R.~Collobert, and J.~Weston,
Curriculum Learning,
in \emph{Proc. ICML}, 2009.

\end{thebibliography}

\end{document}
